\honest{Your Price for Joining}{Eliezer Yudkowsky}{2009}

In the Ultimatum Game, one player proposes how to split \$10; if the other rejects the split, neither gets anything. Conventional causal decision theory says the responder should take any positive amount, so a penny will do. I am no fan of that theory, and those of us who care about fairness ``may also not accept an offer of one penny.''

Nor do real people. I say that ``most Ultimatum `deciders' offer an even split; and most Ultimatum `accepters' reject any offer less than 20\%.'' In Indonesia, offers of \$30 in a \$100 game were turned down, although that is two weeks' wages. \nb{Yudkowsky's figures match the Wikipedia article the post links, which said at the time that low offers are ``often'' rejected. A meta-analysis of 37 papers puts the average offer at 40\% and the average rejection rate at 16\%.} This, I say, ``is just the way people feel about Ultimatum Games.'' \nb{A study of fifteen small-scale societies found that the self-interest prediction fails in every one, and that play varies much more between societies than earlier studies had shown.}

Joining a group is similar, I suggest. ``Has it been studied? I'd hope so...'' \nb{A related question had been: Albert Hirschman's \textsc{Exit, Voice, and Loyalty} (1970) is about members of an organization who, when it disappoints them, either leave or press for change from inside.} Say a project to help mugging victims invests its money badly. Joining would still do more good than fighting crime alone, but if you join too easily, why would they change?

``It seems to me,'' I say, that people in the ``atheist\slash{}libertarian\slash{}technophile/sf-fan/etcetera cluster'' set their joining prices ``way way way too high,'' like fifty players each demanding 20\% of the pot. It is ``almost certainly a matter of evolutionary psychology'': our instincts were set for bands of about 40 people who all knew each other. \nb{The only support given is that such situations would often have arisen among hunter-gatherers.} So we underestimate how hard it is to change a large group or a complex plan. Weakness of will and a wish for status may play a part too. I admit I write as an organizer: ``Perhaps I am a little prejudiced.''

My rule of thumb: if joining would still do net good, and the issue is not worth your fixing it yourself, do not withhold your support. If it is worth that much, join and fix it. If the group will not let a competent person fix it and no one else is harmed, ``it may be time for a little blackmail.'' Is this extreme? ``Oh, maybe.'' A project's leaders should have some reason to answer to their supporters.

But usually, I observe, people set their prices too high. It gets tiring to read ``the 451st public email'' saying the project is not worth joining until the website has a sans-serif font. We ``probably'' want a group norm: if the issue is not worth fixing yourself and is not bad faith, it is not worth refusing to contribute to a cause you think worthwhile.
